\documentclass[12pt]{article}
\usepackage[T1]{fontenc}
\usepackage[utf8]{inputenc}
\usepackage{lmodern}
\usepackage{textcomp}
\usepackage{enumitem}
\usepackage{geometry}
\geometry{margin=1in}
\begin{document}
\begin{center}
Answer Key - Constitutional Law - Undergraduate Level Assessment 9 \
\small (Answers are ordered exactly as in the question paper.)
\end{center}

\section*{Multiple Choice}
\begin{enumerate}[label=\arabic*.]
\item \textbf{Answer:} C
\\ \textit{Options:} A) Domestic law always prevails over international law; B) Only customary international law prevails over domestic law; C) Obligations under international law prevail over domestic law; D) Constitutional obligations always prevail over obligations under international law
\\ \textit{Note:} Obligations under international law prevail over domestic law because of the principle of pacta sunt servanda, which asserts that treaties and international agreements must be observed, regardless of conflicting domestic legislation.
\item \textbf{Answer:} D
\\ \textit{Options:} A) The ultimate source of a legal system's morality.; B) The rule that distinguishes norms from habits of obedience.; C) The constitution of a state.; D) A presupposition that facilitates our understanding of the legal system.
\\ \textit{Note:} Kelsen's Grundnorm serves as a fundamental presupposition that underpins the validity of a legal system, allowing for a coherent understanding of the hierarchy and structure of laws within that system.
\item \textbf{Answer:} B
\\ \textit{Options:} A) Piracy jure gentium is subject to flag State jurisdiction; B) Piracy jure gentium is subject to universal jurisdiction; C) Piracy jure gentium is subject to port State jurisdiction; D) Piracy jure gentium is subject to nationality-based jurisdiction
\\ \textit{Note:} Piracy jure gentium is considered a crime against all nations, allowing any state to exercise universal jurisdiction over pirates, regardless of where the crime occurred or the nationality of the perpetrators or victims.
\item \textbf{Answer:} A
\\ \textit{Options:} A) He denies that law develops independently of social and economic forces.; B) He claims that law is economically immoral.; C) He rejects a positivist account of law.; D) He opposes any sociological analysis of the law.
\\ \textit{Note:} Posner argues that law is heavily influenced by social and economic forces, suggesting that it cannot develop in isolation from these contextual factors, thereby denying its autonomy.
\item \textbf{Answer:} B
\\ \textit{Options:} A) John Rawls; B) Stammler; C) Jerome Hall; D) John Finns
\\ \textit{Note:} Stammler is recognized for his contribution to the theory of natural law by asserting that its principles are not fixed but can evolve in response to changing societal conditions and values.
\end{enumerate}

\section*{True / False}
\begin{enumerate}[label=\arabic*.]
\setcounter{enumi}{5}
\item \textbf{Answer:} False
\\ \textit{Options:} A) True; B) False
\\ \textit{Note:} The statement is false because Gemeinschaft refers to community-oriented societies characterized by close personal relationships and traditional values, while Gesellschaft describes more impersonal, urban societies that may or may not embrace democratic principles.
\item \textbf{Answer:} True
\\ \textit{Options:} A) True; B) False
\\ \textit{Note:} In Hohfeld's scheme, power and liability are oppositional jural concepts where the exercise of power by one party creates a corresponding liability for another, thus establishing a contradiction between these two relationships.
\item \textbf{Answer:} True
\\ \textit{Options:} A) True; B) False
\\ \textit{Note:} This statement is true because Salmond's definition emphasizes that law consists of established principles and rules that are acknowledged and enforced by the state to ensure justice within society.
\item \textbf{Answer:} False
\\ \textit{Options:} A) True; B) False
\\ \textit{Note:} Dworkin's theory of equality of resources encompasses a broader view that includes the distribution of resources, opportunities, and individual preferences, rather than being limited to the amount of income tax paid.
\item \textbf{Answer:} False
\\ \textit{Options:} A) True; B) False
\\ \textit{Note:} The rule of recognition applies to both unitary and federal systems as it establishes the criteria for identifying valid legal norms, regardless of the structure of the constitution.
\end{enumerate}

\section*{Long Form Response}
\begin{enumerate}[label=\arabic*.]
\setcounter{enumi}{10}
\item \textbf{Answer:} The UK Constitution is unique in its uncodified status, meaning there is no single, written document that serves as the foundational legal framework for the country. Instead, it is derived from a variety of sources, including statutes, common law, conventions, and works of authority. This structure allows for a flexible and adaptable legal system that can evolve over time, reflecting changes in society and governance. However, the lack of a definitive constitutional text can lead to ambiguities and debates over the interpretation of laws, which impacts the stability and predictability of governance. Consequently, while the uncodified nature of the UK Constitution allows for responsiveness to change, it also presents challenges in maintaining clear legal standards and accountability.
\\ \textit{Note:} The answer correctly identifies the UK Constitution as uncodified, detailing its diverse sources-statutes, common law, conventions, and authoritative texts-which together create a flexible legal framework that can adapt to societal changes while also acknowledging the potential for ambiguity due to the absence of a single constitutional document.
\item \textbf{Answer:} Emile Durkheim's concepts of mechanical and organic solidarity illustrate how societal cohesion influences the nature of law within a community. In societies characterized by mechanical solidarity, where individuals share similar values and beliefs, the law tends to be repressive; it imposes conformity and punishes deviations to maintain social order. Conversely, in societies marked by organic solidarity, which arises from the interdependence of diverse individuals with specialized roles, the law becomes restitutive, focusing on restoring balance and facilitating cooperation rather than punishment. This distinction underscores how the structure of society can shape legal frameworks, reflecting the underlying social bonds that unify individuals.
\\ \textit{Note:} correct because it effectively links Durkheim's distinction between mechanical and organic solidarity to the corresponding legal frameworks, demonstrating how societal cohesion shapes the repressive nature of laws in homogeneous societies versus the more restitutive laws in heterogeneous, interdependent communities.
\end{enumerate}

\vfill
\noindent\textit{End of Answer Key}
\end{document}